% !TeX root = ../papers/nc-p3.tex
\section{両側一般の高階数窓の線型復元と八固有直線による全同型}
\label{sec:marked-mixed-bracket-reconstruction}
\begin{theorem}\label{thm:marked-mixed-bracket-reconstruction}
一点を指定した窓で、二つの三平面が同じ空間を張るという仮定を外す。
正規化座標の交叉交換子を
\[
  C_{bd,ac}=[Q_b,P_d]_{ac}=-[G_a,F_c]_{bd},
  \qquad C\in\mathfrak{sl}(V)\otimes\mathfrak{sl}(U)
\]
とする。任意の $C$ に対して印付き二切断方程式の全解は、下記二つの
線型写像の各アフィン逆像と階数一テンソル錐との交叉で決まる。
分裂混合点を通る十五次元成分の一般の窓には、両方の写像が単射となる印がある。
その印では $C$ と $C$ の転置配置のそれぞれから六行列が有理演算だけで逆倍率まで復元できる。
従ってこの開部分の正規化印付き表示は、実際の $C$ の三十二次元軌跡の上で
二つの逆倍率を選ぶものとなる。一般点は $g=q=6$、反対階数16である。

この一般点で $C$ を跡対合により $\mathfrak{sl}(V)\to\mathfrak{sl}(U)$ とすれば
可逆で、$C^*C$ は相異なる八固有値を持つ。この条件と両線型写像の単射性を満たす全ての窓では、
後述の有限個の数表の一致が印を保つ添字保持同型の必要十分条件である。
印を忘れた同型では、相手の有限な印について同じ比較を行えばよい。
\end{theorem}

\begin{proof}[混合交換子からの二つの因子]
まず $C_{ij}$ を一方の三行三列の跡零行列として読む。第一の線型空間を
\[
 L_C=\{X\in\mathfrak{sl}_3^3:
       \operatorname{tr}(X_i C_{kj}+X_k C_{ij})=0\ (i\le k,\ 1\le j\le3)\}
\]
とし、第二を $R_C=L_{(C_{ji})}$ とする。$C_{ij}=[X_i,Y_j]$ なら跡の巡回性で
$X\in L_C,Y\in R_C$ である。係数のブラケットにより
\[
  m_C:L_C\otimes R_C\longrightarrow (\mathfrak{sl}_3)^9,\qquad
                  X\otimes Y\longmapsto ([X_i,Y_j])_{ij}
\]
を作る。この線型写像の $C$ 上の逆像の階数一テンソルは全て、かつそれだけが
求める組である。非零の各テンソルは一つの非零成分を選んで行と列を読むと有理的に分解でき、
自由度は $(X,Y)\mapsto(hX,h^{-1}Y)$ となる。$C=0$ の場合の零テンソルには
一方が零の全組を加える（有限窓の独立な三列ではそもそも不要である）。
他方の組にも $-C_{ij,kl}$ の二添字対を入れ替えて同じ判定を行う。
これで任意の例外次元にも二切断方程式そのものの必要十分条件を得る。
各写像が単射であれば線型逆像は空か一点なので、両方の線型整合性と各一点の
二次小行列式の消失が、実際の $C\ne0$ の明示的条件である。

定理\ref{thm:marked-point-commutators}の残余群 $(A,B,s)$ は $C$ に
$s(\operatorname{Ad}_B\otimes\operatorname{Ad}_A)$ と作用する。
固定した $C$ の二つの逆倍率は $A,B$ 各々のスカラーによって独立に動く。
よって単射な開集合では $[C]$ の $\mathrm{PGL}(V)\times\mathrm{PGL}(U)$ 軌道が、
印を保つ添字保持同型の完全不変量となる。
\end{proof}

\begin{proof}[実在する十五次元成分での単射性と単純固有値]
分裂混合点 $G=D,F=-D,Q=P=K$ で四組の跡零行列の九十六座標を使う。
この点の正規化印付き局所環は滑らかで次元34である。実際、元の二切断テンソル軌跡は
命題\ref{prop:split-mixed-transverse-curve}のSchur微分と二十余法線で滑らかな次元76である。
単純な二十点のうち一つを選ぶ被覆は局所でエタールであり、対応する四点は
各辺の階数3の部分評価から復元する。基底でこれを標準点にする表示へ移る際は
$\mathrm{GL}_4^4$ の枠を用いる。標準点の各直線を保つ枠は四因子で次元52、
二つの反対対合を単位行列にすると18条件、補文字の三つずつの持上げは
跡零化の12条件で一意に決まる。これらは可逆行列の選択・線型方程式で滑らかな
次元22の枠の繊維となる。従って正規化表示の次元は $76+22-64=34$ であり、
九十六変数の方程式の微分階数は62である。
同成分の一般点の $g=q=6$ と反対階数16は既知の横断曲線上の非消失から従う。

素数1009で微分の可逆な六十二次小行列を取り、残りの三十四自由座標を動かす。
この六十二式だけの完備局所環は滑らかで、標数零の全方程式の滑らかな同次元の商と一致する。
従って整数係数の自由座標を $1009$ の倍数だけ動かし、この小行列の単元性で逐次補正すると
実際の標数零の局所族の $1009$ 進点が得られる。より具体的には、形式的陰関数の係数は
$\mathbf Z_{1009}$ に入り、高次係数の収束と残りの全恒等式を保ったまま代入できる。
\path{computations/verify_mixed_bracket_reconstruction.py} は一つの自由方向で72桁まで
九十六変数を求め、全方程式と各段の剰余を検算する。
二つの読み方の各々で左右の十八行の階数はともに18、従って両核とも次元6、
積写像の三十六列は独立である。検算は単元ピボットを選び、残る行が一斉に素数で割れるときは
共通に割って有効桁数を減らす。残りが有効桁まで零となる前に全行にピボットが得られる。
同じ点で、八次の $C$ の行列式の付値は8、$C^*C$ の特性多項式の判別式の付値は62である。
いずれも精度72未満なので真の非消失を証明する。よって各条件はこの印付き成分で稠密開で、
印を忘れた成分の一般点にもそのような印が存在する。二倍率の復元から $C$ の像の
次元は $34-2=32$、射影化して31となる。この検算は個別の閉点の全次数AS判定には使用しない。
\end{proof}

\begin{proof}[有限数表の同型判定]
$C^*C$ は跡対合に関して自己随伴で、八固有直線は互いに直交し各々非退化である。
その順序を取り、最初の固有値が1になるよう $C$ の全体倍率を二通り選ぶ。
各固有直線で $\operatorname{tr}(u_i^2)=1$ となる二つの代表を選び、
$v_i=C u_i$、$\lambda_i=\operatorname{tr}(v_i^2)$ とする。有限表は
\[
    (\lambda_i)_i,\qquad (\operatorname{tr}(u_i u_j u_k))_{ijk},\qquad
                            (\operatorname{tr}(v_i v_j v_k))_{ijk}
\]
である。基底順序、全体倍率、八つの符号を全通り許す。
表の一致から $u_i$ 同士と $v_i$ 同士を送る線型同型は、両方の跡対合を保つ。
さらに単位を単位へ送って延長すると全行列積を保つ。実際、積のスカラー成分は
二元の跡の三分の一、残る成分は三元の順序付き跡を各基底のノルムで割ったものである。
従って全行列環の二つの単位的同型となり、行列単位の像からそれぞれ共役である。
二共役は $C$ の作用を可換にするので上記の完全不変量が一致する。逆は各操作の共変性による。
向きを持つ三元の跡を用いるので転置や添字の交換を余分に同一視しない。
無印の入力には第一の窓で一つの適格な印を固定し、第二の窓の同じ位置の全単純点を調べる。
同型は印と全ての階数条件を運び、逆にどれかでの印付き同型は無印の同型を与える。
\end{proof}
